\documentclass[12pt]{article}
\usepackage[utf8]{inputenc}
\usepackage[T1]{fontenc}
\usepackage{hyperref}
\title{Savara Language Resources for K-12 Curriculum Development}
\author{Compiled by Claude Sonnet 5.5 (claude-sonnet-5-5) for Andr\'as Kornai}
\date{}
\begin{document}
\maketitle

\section{Introduction}
The code svr (Glottocode sava1244) is a retired entry. Glottolog lists it as ``Savara'' under bookkeeping, with the reason ``non-existent'' (change request 2015-010, effective 2016-01-15). That also explains why the family is given as ``Bookkeeping''. ``Savara'' is an alternate name of Sora, which has its own code (srb, Glottocode sora1254) and is Austroasiatic (Munda). Resources for Sora belong in the srb dossier, so I have not listed any here, to avoid mixing up the two entries.

\section{Primary Educational Resources}
None are listed here, for the reason given above.

\section{Linguistic and Reference Materials}
\begin{itemize}
\item Glottolog record for ISO code svr, showing the retirement: \url{https://glottolog.org/resource/languoid/iso/svr}
\end{itemize}

\section{Religious and Community Resources}
None are listed here.

\section{Digital and Audio Resources}
None are listed here.

\section{Curriculum Development Context}
There is no language to develop a curriculum for under this code. Curriculum work should be recorded under Sora (srb).

\end{document}